\begin{abstract}
Multi-agent systems (MAS) built from language models are increasingly deployed on the premise that collaboration among agents outperforms a single-agent system (SAS), yet when collaboration pays remains poorly understood. The most comprehensive study to date answers with an empirical capability-saturation rule, offered as a practical selection rule rather than a universal law: coordination helps while the underlying SAS is weak and hurts once it is strong, with the crossing near 45\% single-agent accuracy. That rule was derived on autonomous benchmarks, where no user takes part. We test it where agents increasingly operate: in conversation with a user under a written policy, on $\tau^2$-bench, the dual-control benchmark, where the user supplies information turn by turn and, in its telecom domain, acts on the shared environment as well. On its airline, retail and telecom domains we compare a SAS with a centralized MAS built from the same model (an orchestrator, a database agent and a policy advisor) for six models from three families (four, from two families, in the primary analysis), over about 7,950 simulated conversations, with each system held to one completion-token ledger per episode. The study was pre-registered with a public change log; the MAS is reported as built, which leaves one pre-registered validity bar unmet and labelled as a deviation (its planner never saw the tool schemas). Under a fixed simulated user no cell shows a MAS benefit: the fitted gain is negative across the observed capability range, with a weighted mean of $-28$ points (95\% CI $[-44, -12]$) over the 12 cells of the primary fit, and the MAS is worse in 12 of the 13 cells that meet the mechanical validity bars and tied in the thirteenth, a cell near 45\% where the benchmark's default simulator yields a 10-point MAS gain. On the two cells where the principle predicts a benefit, the MAS loses by 13 and 22 points, and removing the budget ceiling brings it at best level with the SAS at 5.5 times its median token spend. The design could not resolve whether the gain declines with capability: a capability index measured on held-out domains gives a slope of $-0.41$ with a 95\% CI of $[-1.51, +0.68]$, because model-by-domain effects dominate the variation between cells. The mechanism that would justify a MAS here is not detectable: the MAS shows no measurable reduction in wrongful database writes and converts other failures into budget exhaustion and inaction. For agents that serve people, these results make the SAS the default for this architecture and leave multi-agent designs to show compute-matched evidence from the setting itself. We release all trajectories, budget ledgers, prompts and code.
\end{abstract}
